\section{Sources of Bias}

\begin{frame}{First: ``Bias'' Means Three Different Things}

    \begin{columns}[T]
        \begin{column}{0.32\textwidth}
            \begin{tcolorbox}[colback=gray!8, colframe=gray!60, title=Statistical bias]
                \small
                $\mathbb{E}[\hat{\theta}] - \theta \neq 0$.

                \vspace{0.4em}
                The estimator is systematically off. Neutral term; half of the
                bias--variance trade-off.
            \end{tcolorbox}
        \end{column}
        \begin{column}{0.32\textwidth}
            \begin{tcolorbox}[colback=raiBlue!5, colframe=raiBlue!60, title=Inductive bias]
                \small
                The assumptions a model class makes.

                \vspace{0.4em}
                Convolutions assume locality; linear models assume additivity.
                \textbf{Necessary} --- without it there is no generalisation.
            \end{tcolorbox}
        \end{column}
        \begin{column}{0.32\textwidth}
            \begin{tcolorbox}[colback=raiRed!5, colframe=raiRed!70, title=Societal / harmful bias]
                \small
                Systematically different quality of service or outcome across groups of people.

                \vspace{0.4em}
                \textbf{This lecture's subject.}
            \end{tcolorbox}
        \end{column}
    \end{columns}

    \vspace{1.2em}

    \begin{block}{}
        \centering
        The three are unrelated. A model can be statistically unbiased and still harmful; a strong
        inductive bias can either mask or expose a societal one.
    \end{block}
\end{frame}

\begin{frame}{Bias Enters at Every Stage, Not Just ``The Data''}

    \centering
    \resizebox{\textwidth}{!}{%
    \begin{tikzpicture}[
        node distance=0.55cm,
        stage/.style={rectangle, draw, rounded corners, minimum height=0.95cm, text width=2.05cm,
                      align=center, font=\scriptsize\bfseries, fill=raiBlue!10},
        tag/.style={font=\scriptsize, align=center, text width=2.2cm, text=raiRed!80!black},
        arr/.style={-{Latex[length=2mm]}, thick}
    ]
        \node[stage] (world) {The world};
        \node[stage, right=of world] (collect) {Collect \& sample};
        \node[stage, right=of collect] (label) {Define target \& label};
        \node[stage, right=of label] (train) {Train};
        \node[stage, right=of train] (eval) {Evaluate};
        \node[stage, right=of eval] (deploy) {Deploy};

        \foreach \a/\b in {world/collect, collect/label, label/train, train/eval, eval/deploy}
            \draw[arr] (\a) -- (\b);

        \node[tag, below=0.35cm of world]   {Historical bias};
        \node[tag, below=0.35cm of collect] {Representation bias};
        \node[tag, below=0.35cm of label]   {Measurement bias};
        \node[tag, below=0.35cm of train]   {Aggregation \& learning bias};
        \node[tag, below=0.35cm of eval]    {Evaluation bias};
        \node[tag, below=0.35cm of deploy]  {Deployment bias};

        \draw[arr, dashed, raiRed] (deploy.north) .. controls +(0,1.2) and +(0,1.2) .. (world.north)
            node[midway, above, font=\scriptsize\itshape, text=raiRed] {feedback loop};
    \end{tikzpicture}}

    \vspace{1em}
    {\scriptsize Taxonomy after H.~Suresh and J.~V.~Guttag, ``A Framework for Understanding Sources of
    Harm throughout the Machine Learning Life Cycle,'' EAAMO 2021,
    \href{https://arxiv.org/abs/1901.10002v5}{arXiv:1901.10002v5}.}
\end{frame}

\begin{frame}{Historical and Representation Bias}

    \begin{columns}[T]
        \begin{column}{0.48\textwidth}
            \textbf{Historical bias}
            \vspace{0.3em}

            Present even with perfect sampling and perfect measurement, because the world the data
            records is itself unequal.

            \vspace{0.6em}
            \begin{itemize}\small
                \setlength\itemsep{0.35em}
                \item An image search for ``CEO'' that reflects who currently holds the job.
                \item A loan model trained on approvals from an era of discriminatory lending.
            \end{itemize}

            \vspace{0.6em}
            \begin{tcolorbox}[colback=gray!8, boxrule=0pt]
                \small More data does not fix this. It is a question about the objective, not
                the sample.
            \end{tcolorbox}
        \end{column}
        \begin{column}{0.48\textwidth}
            \textbf{Representation bias}
            \vspace{0.3em}

            The development sample under-represents part of the use population, so the model
            generalises poorly there.

            \vspace{0.6em}
            \begin{itemize}\small
                \setlength\itemsep{0.35em}
                \item A dermatology model trained mostly on light skin.
                \item A speech model trained on one accent.
                \item Gender Shades: existing face benchmarks were 79.6\% and 86.2\%
                      lighter-skinned subjects.
            \end{itemize}

            \vspace{0.6em}
            \begin{tcolorbox}[colback=gray!8, boxrule=0pt]
                \small More \emph{targeted} data does help --- but you must first know the
                composition of your data.
            \end{tcolorbox}
        \end{column}
    \end{columns}
\end{frame}

\begin{frame}{Measurement Bias: The Label Is a Proxy}

    You almost never observe the thing you care about. You observe something correlated with it.

    \vspace{0.8em}
    \centering
    \small
    \begin{tabular}{p{3.9cm} p{4.0cm} p{4.0cm}}
        \toprule
        \textbf{What you want} & \textbf{What you record} & \textbf{What the gap encodes} \\
        \midrule
        Health need & Health-care spending & Unequal access to care \\
        \addlinespace
        Committing a crime & Being arrested & Unequal policing intensity \\
        \addlinespace
        Job performance & Was hired / was promoted & Past hiring decisions \\
        \addlinespace
        Creditworthiness & Repaid a loan they were \emph{granted} & Who was offered credit \\
        \addlinespace
        Content quality & Clicks and watch time & Engagement, not quality \\
        \bottomrule
    \end{tabular}

    \vspace{1em}
    \raggedright
    \begin{block}{The one question to ask on every project}
        \centering
        \textit{``My label is X. What exactly had to happen in the world for X to be recorded,
        and does that differ by group?''}
    \end{block}
\end{frame}

\begin{frame}{Aggregation, Learning, Evaluation and Deployment Bias}

    \begin{itemize}
        \setlength\itemsep{0.7em}
        \item \textbf{Aggregation bias.} One model for populations that genuinely behave
              differently. Suresh and Guttag point to medicine: patients of different sexes with
              similar underlying conditions can present and progress differently, so one model is
              optimal for neither.

        \item \textbf{Learning bias.} The modelling choice itself amplifies a disparity.
              Optimising average loss lets a large majority group dominate the gradient;
              aggressive compression or regularisation tends to cost the tails most.

        \item \textbf{Evaluation bias.} The benchmark does not represent the use population,
              so the reported number is not the number you will get. This is exactly the
              Gender Shades finding about existing face benchmarks.

        \item \textbf{Deployment bias.} The model is used differently from how it was framed.
              A risk score built to \emph{rank} people for services becomes a
              \emph{decision} about their liberty; a triage aid becomes an autopilot.
    \end{itemize}

    \vspace{0.4em}
    \centering
    \textit{Four of the seven sources arise \emph{after} the data has been collected.}
\end{frame}

\begin{frame}{Feedback Loops: The Model Changes Its Own Training Data}

    \centering
    \begin{tikzpicture}[
        node distance=1.7cm and 2.4cm,
        box/.style={rectangle, draw, rounded corners, minimum height=1.0cm, text width=3.0cm,
                    align=center, font=\small},
        arr/.style={-{Latex[length=2.5mm]}, thick}
    ]
        \node[box, fill=raiBlue!12] (model)  {Model predicts high risk in area A};
        \node[box, fill=gray!12, right=of model] (action) {More patrols sent to area A};
        \node[box, fill=raiRed!12, below=of action] (record) {More arrests recorded in area A};
        \node[box, fill=raiRed!20, below=of model]  (retrain) {Retraining confirms area A is high risk};

        \draw[arr] (model)  -- (action);
        \draw[arr] (action) -- (record);
        \draw[arr] (record) -- (retrain);
        \draw[arr] (retrain) -- (model);
    \end{tikzpicture}

    \vspace{0.8em}

    \begin{itemize}\small
        \setlength\itemsep{0.35em}
        \item The loop is self-confirming: the prediction generates the evidence for itself.
        \item Same shape in recommenders (you only observe clicks on what you showed),
              in credit (you only observe repayment for approved applicants), and in hiring.
        \item \textbf{Detection:} hold out a slice from the policy, or log what the model
              \emph{suppressed}, not just what it surfaced.
    \end{itemize}
\end{frame}

\begin{frame}{``Just Drop the Sensitive Attribute'' Does Not Work}

    \textbf{Fairness through unawareness:} remove race, gender, age from the feature set and
    declare the model neutral.

    \vspace{0.8em}

    \begin{columns}[T]
        \begin{column}{0.55\textwidth}
            \begin{itemize}\small
                \setlength\itemsep{0.4em}
                \item Real features are \textbf{correlated} with the removed attribute: postcode,
                      first name, school, device type, browsing history, purchase categories.
                \item A capable model \emph{reconstructs} the attribute from the proxies. Removing
                      the column removes your ability to \textbf{measure} the disparity, not the
                      disparity itself.
                \item The Amazon reporting is exactly this: no gender field, gendered signal.
                \item Corollary: you often need the sensitive attribute --- held securely and used
                      only for auditing --- in order to check for harm.
            \end{itemize}
        \end{column}
        \begin{column}{0.43\textwidth}
            \centering
            \begin{tikzpicture}[scale=0.9,
                n/.style={circle, draw, minimum size=0.85cm, font=\scriptsize, align=center},
                arr/.style={-{Latex[length=2mm]}, thick}]
                \node[n, fill=raiRed!15] (g)  at (0,2)     {$A$};
                \node[n, fill=gray!12]   (z1) at (-1.6,0.6) {zip};
                \node[n, fill=gray!12]   (z2) at (0,0.6)    {name};
                \node[n, fill=gray!12]   (z3) at (1.6,0.6)  {club};
                \node[n, fill=raiBlue!15](y)  at (0,-1.1)   {$\hat{Y}$};
                \draw[arr] (g) -- (z1); \draw[arr] (g) -- (z2); \draw[arr] (g) -- (z3);
                \draw[arr] (z1) -- (y); \draw[arr] (z2) -- (y); \draw[arr] (z3) -- (y);
                \draw[arr, dashed, raiRed, thick] (g) to[bend right=60] node[left, font=\scriptsize]
                    {removed} (y);
            \end{tikzpicture}
            \vspace{0.2em}
            {\scriptsize Cutting the direct edge leaves every indirect path intact.}
        \end{column}
    \end{columns}
\end{frame}

\begin{frame}{Bias Lives in Learned Representations Too}

    Word embeddings trained on news text place words in a space where a
    \textbf{gender direction} is recoverable, and gender-neutral occupation words project onto it.

    \vspace{0.12em}
    \centering
    \begin{tikzpicture}[scale=0.61, font=\scriptsize]
        \draw[{Latex[length=2mm]}-{Latex[length=2mm]}, thick, raiRed] (-4.7,0) -- (4.7,0);
        \node[font=\scriptsize\bfseries, text=raiRed, anchor=east] at (-4.85,0) {she};
        \node[font=\scriptsize\bfseries, text=raiRed, anchor=west] at ( 4.85,0) {he};
        \node[font=\scriptsize\itshape, text=raiRed] at (0,-0.95)
            {gender direction $\approx \overrightarrow{\text{he}} - \overrightarrow{\text{she}}$};

        \foreach \x/\w in {-4.2/homemaker, -2.8/nurse, -1.4/receptionist,
                            1.4/financier, 2.8/architect, 4.2/philosopher} {
            \fill[raiBlue] (\x,0) circle (2.2pt);
            \node[text=raiBlue!80!black, font=\tiny, rotate=35, anchor=west, xshift=3pt] at (\x,0.06) {\w};
        }
        \fill[gray] (0,0) circle (1.8pt);
        \node[text=gray, anchor=north] at (0,-0.15) {neutral};
    \end{tikzpicture}

    \vspace{0.27em}

    \begin{itemize}\scriptsize
        \setlength\itemsep{0.1em}
        \item Bolukbasi et al.\ showed the analogy solver completes
              \emph{man\,:\,computer programmer\,::\,woman\,:\,X} with ``homemaker,''
              and proposed a projection-based debiasing method.
        \item \textbf{Caveat worth teaching:} later work (Gonen \& Goldberg, 2019) showed such
              projection mainly \emph{hides} the bias --- debiased vectors still cluster by
              stereotype, so the association is recoverable.
        \item Practical reading: a fix that makes a metric go down has not necessarily made the
              system fair.
    \end{itemize}

    \vspace{0.09em}
    {\scriptsize Schematic drawn for this slide; the gender direction, the analogy result and the
    six occupation words (from the paper's extreme ``she''/``he'' lists) are from Bolukbasi, Chang,
    Zou, Saligrama \& Kalai, ``Man is to Computer Programmer as Woman is to Homemaker? Debiasing
    Word Embeddings,'' NeurIPS 2016, \href{https://arxiv.org/abs/1607.06520v1}{arXiv:1607.06520v1},
    Figure 1. Follow-up: Gonen \& Goldberg, ``Lipstick on a Pig,'' NAACL 2019,
    \href{https://arxiv.org/abs/1903.03862v2}{arXiv:1903.03862v2}.}
\end{frame}
